% Dokumentace projektu IFJ26 / IAL – šablona (pouze struktura)
% Překlad: pdflatex dokumentace.tex (2x kvůli odkazům)
\documentclass[a4paper,11pt]{article}

\usepackage[utf8]{inputenc}
\usepackage[T1]{fontenc}
\usepackage[czech]{babel}
\usepackage[margin=2.5cm]{geometry}
\usepackage{lmodern}
\usepackage{microtype}
\usepackage{booktabs}
\usepackage{tabularx}
\usepackage{array}
\usepackage{enumitem}
\usepackage{relsize}
\usepackage{graphicx}
\usepackage[justification=centering,singlelinecheck=false]{caption}
\usepackage{pdflscape}
\usepackage{xcolor}
\usepackage{tikz}
\usepackage{circuitikz}
\usetikzlibrary{automata,positioning,arrows.meta,shapes.geometric,fit}
\usepackage[hidelinks]{hyperref}
\usepackage{amssymb,booktabs}
\usepackage{float}

\setlist{itemsep=2pt,topsep=4pt}

% Styl a značení pro diagram konečného automatu (fsm-diagram.tex)
\tikzset{
  ka/.style={
    ->, >={Stealth[length=2mm]}, auto, semithick,
    font=\footnotesize\ttfamily,
    node distance=2cm,
    initial text={},
    every state/.style={
      draw, rectangle, rounded corners=2mm, align=center,
      minimum height=7mm, minimum width=9mm, inner sep=2pt, fill=white,
    },
    every loop/.style={min distance=6mm, looseness=7},
  },
}
% Popisky přechodů
\newcommand{\kaLF}{{\normalfont\textit{LF}}}
\newcommand{\kaTAB}{{\normalfont\textit{TAB}}}
\newcommand{\kaSP}{\textvisiblespace}
\newcommand{\kaQ}{\textquotedbl}
\newcommand{\kaBS}{\textbackslash}
% Podtržítko v rámečku -- odlišitelné od mínusu
\newcommand{\kaUS}{{\setlength{\fboxsep}{0.6pt}\setlength{\fboxrule}{0.3pt}\fbox{\vphantom{Ag}\_}}}

% Zástupný text – před odevzdáním všechny výskyty odstranit
\newcommand{\TODO}[1]{\textcolor{red}{[#1]}}

% Údaje o týmu (vyplnit)
\newcommand{\varianta}{vv-BVS}
\newcommand{\rozsireni}{FUNEXP, CONTROL, BOOLTHEN, LISTS}

\begin{document}

%------------------------------------------------------------
% 1. strana
%------------------------------------------------------------
\begin{titlepage}
  \centering
  {\Large Vysoké učení technické v~Brně\\
   Fakulta informačních technologií\par}
  \vspace{2cm}
  {\Huge\bfseries Implementace překladače\\imperativního jazyka IFJ26\par}
  \vspace{0.5cm}
  {\large Dokumentace k~projektu z~předmětů IFJ a~IAL\par}
  \vspace{1cm}
  {\large Tým xnovakm00, varianta \varianta\par}
  \vspace{1.5cm}
  \begin{center}
    \begin{tabular}{@{}l l l r@{}}
        \toprule
        \textbf{Jméno a příjmení} & \textbf{Login} & \textbf{Osobní číslo} & \textbf{Body [\%]}\\
        \midrule
        \textbf{Milan Novák} & \textbf{xnovakm00} & {281702} & 25\\
        Šimon Hrouda & xhrouds00 & 281730 & 25\\
        Vojtěch Ondráček & xondrav00 & 282118 & 25\\
        Lukáš Petříček & xpetril00 & 281978 & 25\\
        \bottomrule
    \end{tabular}
  \end{center}

  \vspace{1cm}
  \begin{tabularx}{\textwidth}{@{}l X@{}}
    \textbf{Varianta:} & \varianta\\[4pt]
    \textbf{Implementovaná rozšíření:} & \rozsireni\\[4pt]
    \textbf{Git repozitář:} & \href{https://github.com/phuid/ifj26-compiler}{https://github.com/phuid/ifj26-compiler}\\
  \end{tabularx}

  \vfill
  {\large \today}
\end{titlepage}

\tableofcontents
\newpage

%------------------------------------------------------------
\section{Rozdělení práce v týmu}
\label{sec:prace}
Milan Novák: Výškově vyvážený strom, generátor cílového kódu\\
Šimon Hrouda: Preprocesor, rekuzivní sestup, precedenční analýza výrazu\\
Vojtěch Ondráček: Lexikální analyzátor, rekuzivní sestup, semantická analýza\\
Lukáš Petříček: Testy, semantická analýza

\subsection{Dělení bodů}
Všichni členové týmu spolupracovali a kooperovali adekvátně, nemáme důvod pro nestandardní rozdělení bodů.

%------------------------------------------------------------
\section{Spolupráce a kvalita kódu}
\label{sec:spoluprace}

\subsection{Způsob spolupráce}
Týmová komunikace probíhala přes aplikaci Discord. Každý týden se celý tým účastnil 2-3 hodinové porady. Překladač jsme designovali primárně kooperativní metodou během týmových porad. Často jsme využívali metody párového programování, buď fyzicky na FITu, nebo online přes    Discord.
% Komunikační kanály, schůzky, vedení repozitáře.

\subsection{Code review a schvalování PR}
Každý pracoval výhradně ve vlastní větvi. Před merge do hlavní větve vždy proběhl code review a oprava poznámek. Code review byl vynucen pravidlem repozitáře. V některých případech code review probíhal přímo jako součást párového programování, což bylo rychlejší než klasický postup implementace -> review -> implementace...

\subsection{Formy testování}
% Max. jednotky vět či odrážek (unit testy, end-to-end testy apod.).
\begin{itemize}
  \item \textbf{Jednotkové testy} -- každou část překladače (lexer, parser,
        výrazy, sémantiku, tabulku symbolů, strom~\dots) jsme testovali zvlášť
        pomocí knihovny GoogleTest. Celkem jde zhruba o~470 testů.
  \item \textbf{End-to-end testy} -- přes 200 krátkých programů v~IFJ26,
        u~kterých víme, jak mají dopadnout. Skript každý program přeloží
        a spustí v~interpretu. Pak zkontroluje, jestli sedí výstup i návratový
        kód. U~chybných programů kontroluje, že překladač vrátí správný
        chybový kód.
  \item \textbf{Automatické kontroly} -- po každém pushi na GitHub se
        automaticky zkontroluje formátování kódu a proběhne statická analýza
        (\texttt{cppcheck}).
\end{itemize}

\subsection{Styl formátování kódu}
% Zvolený styl clang-format (hodnota --style).
Automatické formátování nástrojem clang-format při každém commitu.\\
Byl zvolen výchozí styl LLVM.

%------------------------------------------------------------
\section{Architektura překladače}
\label{sec:architektura}
\subsection{Blokové schéma}
% Moduly + datové toky (tokeny, AST, mezikód, ...).
\begin{figure}[H]
  \centering
  % \begin{tikzpicture}[...] ... \end{tikzpicture}
  % \fbox{\parbox{0.9\textwidth}{\centering\vspace{4cm}
  
    \centering
  \resizebox{0.6\linewidth}{!}{%
    \begin{circuitikz}
      \tikzstyle{every node}=[font=\fontsize{14.4pt}{18.7pt}\selectfont]
      \draw [fill={rgb,255:red,255; green,255; blue,255}, fill opacity=1] (11.875,-5.625) rectangle (22.375,-15.75);
      \draw [fill={rgb,255:red,255; green,255; blue,255}, fill opacity=1] (10.875,-6.25) rectangle (12.875,-7.625);
      \node[anchor=center, align=center] at (11.875,-6.9375) {Zdrojový\\ kód};
      \draw [-{Stealth[scale=1.5]}] (10,-6.875) -- (10.875,-6.875);
      \draw [fill={rgb,255:red,255; green,255; blue,255}, fill opacity=1, rounded corners = 7.8] (13.625,-6.375) rectangle (16.875,-7.375);
      \node[anchor=center, align=center] at (15.25,-6.875) {Preprocesor};
      \draw [-{Stealth[scale=1.5]}] (12.875,-6.875) -- (13.625,-6.875);
      \draw [fill={rgb,255:red,255; green,255; blue,255}, fill opacity=1, rounded corners = 7.8] (13.625,-8.75) rectangle (16.875,-11.625);
      \node[anchor=center, align=center] at (15.25,-10.1875) {Lexikální\\ analýza\\ (lexer)};
      \draw [-{Stealth[scale=1.5]}] (15.25,-7.375) -- (15.25,-8.75)node[pos=0.5, fill={rgb,255:red,255; green,255; blue,255}, fill opacity=1, text opacity=1, inner xsep=0.080cm, inner ysep=0.085cm, rounded corners=0.020cm, font=\fontsize{10.95pt}{14.2pt}\selectfont]{upravený zdrojový kód};
      \draw [-{Stealth[scale=1.5]}] (16.875,-11.125) -- (18.75,-11.125);
      \draw [-{Stealth[scale=1.5]}] (18.75,-10.625) -- (16.875,-10.625);
      \draw [fill={rgb,255:red,255; green,255; blue,255}, fill opacity=1, rounded corners = 7.8] (18.75,-8.75) rectangle (21.875,-11.625);
      \node[anchor=center, align=center] at (20.3125,-10.1875) {Syntaktická\\ analýza};
      \draw [fill={rgb,255:red,255; green,255; blue,255}, fill opacity=1, rounded corners = 7.8] (18.75,-12.625) rectangle (21.875,-15.25);
      \node[anchor=center, align=center] at (20.3125,-13.9375) {Semantická\\ analýza};
      \draw [-{Stealth[scale=1.5]}] (20.375,-11.75) -- (20.375,-12.625)node[pos=0.5, left, fill={rgb,255:red,255; green,255; blue,255}, fill opacity=1, text opacity=1, inner xsep=0.080cm, inner ysep=0.085cm, rounded corners=0.020cm, font=\fontsize{10.95pt}{14.2pt}\selectfont]{Program*};
      \draw [-{Stealth[scale=1.5]}] (18.75,-13.875) -- (16.875,-13.875)node[pos=0.5, above, fill={rgb,255:red,255; green,255; blue,255}, fill opacity=1, text opacity=1, inner xsep=0.080cm, inner ysep=0.085cm, rounded corners=0.020cm, font=\fontsize{10.95pt}{14.2pt}\selectfont]{Program*};
      \draw [fill={rgb,255:red,255; green,255; blue,255}, fill opacity=1, rounded corners = 7.8] (13.625,-12.75) rectangle (16.875,-15.25);
      \node[anchor=center, align=center] at (15.25,-14) {Generátor\\ cílového\\ kódu};
      \draw [fill={rgb,255:red,255; green,255; blue,255}, fill opacity=1] (10.875,-13.125) rectangle (12.875,-14.5);
      \node[anchor=center, align=center] at (11.875,-13.8125) {Cílový\\ kód};
      \draw [-{Stealth[scale=1.5]}] (13.625,-13.875) -- (12.875,-13.875);
      \draw [-{Stealth[scale=1.5]}] (10.875,-13.875) -- (10.125,-13.875);
      \node [font=\fontsize{10.95pt}{14.2pt}\selectfont] at (17.75,-11.5) {token/err};
      \node [font=\fontsize{10.95pt}{14.2pt}\selectfont] at (17.75,-10.25) {lexer\_pls};
    \end{circuitikz}
  }
  % }}
  \caption{Architektura překladače.\\
    \footnotesize * Programem je myšlena níže popsaná datová struktura.
  }
  \label{fig:architektura}
\end{figure}

\subsection{Popis komponent a datových toků}

Zdrojový kód nejprve zpracovává preprocesor, který odstraní prolog a komentáře.

Upravený zdrojový kód předá lexeru, který jej převede na tokeny a ty dále zpracovává syntaktická analýza.

Syntaktická analýza metodami rekurzivního sestupu a precedenční SA zdola nahoru se zásobníkem parsuje tokeny ve stejném pořadí v jakém je předá lexer. Při tom plní tabulku funkcí, tabulky symbolů a staví abstraktní syntaktický strom (AST). To vše předá sémantické analýze ve struktuře Program.

Sémantická analýza provádí typové kontroly a z toho vycházející úpravy AST. Upravenou strukturu Program předá generátoru cílového kódu, který rekurzivně prochází AST a generuje cílový kód.

\subsubsection{Datové struktury}
Nejdůležitějším datovým typem pro syntaktickou a sémantickou analýzu a následné generování cílového kódu je tabulka funkcí~\texttt{FunctionTable}, která odpovídá nejvyššímu scope. Každý funkční symbol~\texttt{FunctionSymbol} z tabulky (tj. každá funkce) drží ukazatele na svůj lokální kořenový scope~\texttt{SymbolTable} a tělo funkce jakožto stromová struktura výrazů a příkazů \texttt{ASTNode}. Propojení datových struktur ilustruje Obrázek~\ref{fig:datove-struktury}.

Komunikaci mezi lexikální a sémantickou analýzou zprostředkovává datová struktura \texttt{Token}, která obsahuje typ tokenu a jeho hodnotu (např. konstantu, nebo text identifikátoru).
% Diagram datových struktur překladače (dle src/*.h na větvi main).
% Každý typ je v diagramu právě jednou; rekurzivní struktury (strom scopů,
% AST, uzly BST) jsou znázorněny smyčkou na sebe sama.
% Soubor je samostatný: styly jsou definované přímo v tikzpicture, vyžaduje
% pouze balík tikz a knihovnu arrows.meta (obojí je v preambuli main.tex).
% čárkovaná šipka = vypůjčený ukazatel (bez vlastnictví).

\begin{figure}[H]
  \centering
  \begin{tikzpicture}[
      ds/.style={draw, rounded corners=1mm, fill=black!4, align=left,
                 inner sep=4pt, font=\footnotesize},
      own/.style={-{Stealth[length=2mm]}, semithick},
      ref/.style={-{Stealth[length=2mm]}, semithick, dashed, black!60},
      lab/.style={font=\scriptsize\ttfamily, fill=white, inner sep=1pt},
    ]
    % ---------------------------------------------------------------- uzly
    \node[ds] (prog) at (0,0)
      {\textbf{Program}\\
       \texttt{functions, initialized}};
    \node[ds] (ft) at (0,-1.8)
      {\textbf{FunctionTable} {\scriptsize(dynamické pole)}\\
       \texttt{functions[], function\_count}};
    \node[ds] (fs) at (0,-4.0)
      {\textbf{FunctionSymbol} {\scriptsize(jedna funkce)}\\
       \texttt{name, return\_type, argument\_count}\\
       \texttt{parameters[]}$^{*}$\texttt{, local\_scope, ast\_root}\\
       \texttt{callees[], is\_main, is\_builtin, \dots}};

    \node[ds, anchor=north] (st) at (-3.0,-6.0)
      {\textbf{SymbolTable} {\scriptsize(jeden scope)}\\
       \texttt{parent, children[], bindings}};
    \node[ds] (bst) at (-3.0,-8.6)
      {\textbf{BST} {\scriptsize(výškově vyvážený)}\\
       \texttt{root, compare\_data, \dots}};
    \node[ds] (bnode) at (-3.0,-10.8)
      {\textbf{Node} {\scriptsize(uzel stromu)}\\
       \texttt{value, height, left, right}};
    \node[ds] (sym) at (0.6,-12.9)
      {\textbf{Symbol} {\scriptsize(proměnná / parametr)}\\
       \texttt{name, type, initialized,}\\
       \texttt{constant, owner, assembler\_suffix}};

    \node[ds, anchor=north] (ast) at (3.3,-6.0)
      {\textbf{ASTNode} {\scriptsize(příkaz / výraz)}\\
       \texttt{kind, inferred\_type,}\\
       \texttt{line, column,}\\
       \texttt{data} -- union podle \texttt{kind}:\\
       \texttt{block, if\_statement,}\\
       \texttt{while\_loop, for\_loop,}\\
       \texttt{assignment, call,}\\
       \texttt{binary\_op, unary\_op, \dots}};

    % ---------------------------------------------------------- vlastnictví
    \draw[own] (prog) -- node[lab, right] {functions} (ft);
    \draw[own] (ft) -- node[lab, right] {functions[i] \normalfont\scriptsize(1\,:\,n)} (fs);
    \draw[own] ([xshift=0.9cm]fs.south west) -- node[lab, above left]
      {local\_scope \normalfont\scriptsize(kořenový scope)} (st.north);
    \draw[own] ([xshift=1.7cm]fs.south) -- node[lab, above right]
      {ast\_root \normalfont\scriptsize(kořen: \texttt{AST\_BLOCK})} (ast.north);

    % strom scopů: vnořené scopy jsou tentýž typ SymbolTable
    \draw[own] ([yshift=0.2cm]st.west) .. controls +(150:1.2) and +(210:1.2) ..
      node[lab, left, align=right] {children[i]\\\normalfont\scriptsize(0\,:\,n)}
      ([yshift=0.15cm]st.south west);
    \draw[own] (st) -- node[lab, left] {bindings} (bst);
    \draw[own] (bst) -- node[lab, left] {root} (bnode);
    \draw[own] ([yshift=0.2cm]bnode.west) .. controls +(150:1.2) and +(210:1.2) ..
      node[lab, left] {left, right} ([yshift=0.15cm]bnode.south west);
    \draw[own] (bnode.south) to[out=-90, in=180]
      node[lab, pos=0.4, below left] {value} (sym.west);

    % AST: potomci uzlu jsou opět uzly ASTNode
    \draw[own] ([yshift=0.6cm]ast.east) .. controls +(20:1.6) and +(340:1.6) ..
      node[lab, right, align=left] {\normalfont\scriptsize potomci\\(0\,:\,n)}
      ([yshift=0.6cm]ast.south east);

    % ------------------------------------------------- vypůjčené ukazatele
    \draw[ref] ([yshift=0.15cm]st.south east) .. controls +(330:0.9) and +(30:0.9) ..
      node[lab, right] {parent} ([yshift=0.2cm]st.east);
    \draw[ref] ([xshift=0.9cm]ast.south west) to[out=-90, in=0]
      node[lab, pos=0.45, right] {variable.symbol, target} (sym.east);
    \draw[ref] ([xshift=1.4cm]ast.north) to[out=90, in=0]
      node[lab, pos=0.6, right] {call.function} (fs.east);
  \end{tikzpicture}
  \caption{Diagram datových struktur překladače. Každý typ je uveden jednou,
    smyčka značí rekurzivní strukturu: tabulky symbolů tvoří strom
    zanořených scopů (funkce odkazuje pouze na kořenový scope, vnořené
    scopy jsou jeho potomci), uzly AST mají za potomky opět uzly AST (např.\ příkazy bloku, operandy
    výrazu, větve podmínky nebo tělo cyklu).
    Plná šipka značí vlastnictví (podstruktura je uvolněna spolu
    s~rodičem), čárkovaná šipka vypůjčený ukazatel bez vlastnictví.
    $^{*}$Parametry jsou vypůjčené symboly z~kořenového scope funkce.}
  \label{fig:datove-struktury}
\end{figure}


%------------------------------------------------------------
\section{Lexikální analýza}
\label{sec:lex}
Lexer jsme se rozhodli implementovat jako konečný automat s rozšířením pro rozpoznávání odsazení. Lexikální analyzátor následně pomocí funkce \texttt{lexer\_pls()} vždy vrací další načtený token, případně lexikální chybu. Pokud dojde ke změně odsazení, vrátí speciální token pro zvětšení nebo zmenšení odsazení.

Pokud vyhodnotí, že jde o identifikátor, podívá se ještě do tabulky klíčových slov, která nemohou být identifikátory. Pokud se s některým shoduje, vrátí token tohoto klíčového slova.

Tokeny konstant, které mají hodnotu, přímo konvertuje na příslušnou hodnotu, ať už jde o číslo, nebo řetězec, a vrátí je společně s tokenem. Úvodní prolog a veškeré komentáře, řeší zvlášť preprocesor, který se těchto nepotřebných prvků se zbaví.

Protože lexikální analýza vznikala jako první část implementace projektu, obsahuje i tokeny pro rozšíření, které jsme se nakonec rozhodli nepodporovat.

\subsection{Diagram konečného automatu}
Kvůli přehlednosti je konečný automat lexikálního analyzátoru rozdělen do
šesti částí (obrázky~\ref{fig:ka}--\ref{fig:ka-nl}). Části spolu sdílejí
pouze počáteční stav~\texttt{S}; ostatní stavy a přechody se mezi částmi
neopakují a žádný přechod nevede z~jedné části do druhé. Úplný automat
vznikne sjednocením všech částí, ve kterém se stavy \texttt{S} ztotožní.

Koncové stavy jsou vyznačeny dvojitým okrajem. Pokud automat v~nekoncovém
stavu nemůže pokračovat, jde o~lexikální chybu. Na přechodech se používají
následující značení:
\begin{itemize}
  \item \kaLF{} -- znak konce řádku, \kaTAB{} -- tabulátor, \kaSP{} -- mezera,
  \item \texttt{\kaUS} -- podtržítko (v~rámečku, aby nebylo zaměnitelné se znakem \texttt{-}),
  \item $C = \{c \in \mathrm{ASCII} \mid c \neq \mbox{\kaLF} \wedge c \neq \texttt{\kaQ} \wedge c \neq \texttt{\kaBS}\}$
        -- znaky jednořádkového řetězce bez escape sekvence,
  \item $M = \{c \in \mathrm{ASCII} \mid c \neq \texttt{\kaQ}\}$
        -- znaky víceřádkového řetězce,
  \item \texttt{a-b} -- libovolný znak z~rozsahu \texttt{a} až \texttt{b}.
\end{itemize}

% Diagram konečného automatu lexikálního analyzátoru (dle KA.md).
% Automat je rozdělen na nezávislé části, které sdílejí pouze počáteční
% stav S. Celý automat vznikne sjednocením všech částí (stav S se ztotožní).
% Interní jména uzlů jsou alfanumerická, zobrazené popisky odpovídají KA.md.

\begin{figure}[H]
  \centering
  \begin{tikzpicture}[ka]
    \node[state, initial]   (S)      at (0,0)     {S};
    \node[state, accepting] (Int0)   at (2.4,2.4) {Int0};
    \node[state]            (UDSZ)   at (4.3,4.8) {U\_DoubleStarts\\Zero};
    \node[state]            (DSZ)    at (6.2,2.4) {DoubleStarts\\Zero};
    \node[state, accepting] (Int)    at (2.4,0)   {Int};
    \node[state]            (UInt)   at (2.4,-2)  {U\_Int};
    \node[state, accepting] (Period) at (2.4,-4)  {Period};
    \node[state]            (NumDot) at (6.2,0)   {NumDot};
    \node[state, accepting] (Double) at (8.8,0)   {Double};
    \node[state]            (UDouble) at (8.8,-2.1) {U\_Double};
    \node[state]            (Exp)    at (8.8,2.4) {Exp};
    \node[state]            (SExp)   at (10.3,4.4) {S\_Exp};
    \node[state, accepting] (EDouble) at (11.8,2.4) {E\_Double};
    \node[state]            (EUDouble) at (13.6,4.4) {E\_U\_Double};
    \node[state, accepting] (DFinal) at (11.8,0)  {DoubleFinal};

    \path
      (S)      edge node {0}           (Int0)
               edge node {1-9}         (Int)
               edge node[swap] {.}     (Period)
      % počáteční nula -- smí následovat pouze Double
      (Int0)   edge node {\kaUS}          (UDSZ)
               edge node {0-9}         (DSZ)
               edge node[pos=0.3] {.}  (NumDot)
               edge[out=90, in=90, looseness=2.2] node[pos=0.5] {e,E} (Exp)
      (DSZ)    edge[loop above] node {0-9} ()
               edge[bend left=15] node {\kaUS} (UDSZ)
               edge node {.}           (NumDot)
               edge node {e,E}         (Exp)
      (UDSZ)   edge[loop right] node {\kaUS} ()
               edge[bend left=15] node {0-9} (DSZ)
      % celé číslo
      (Int)    edge[loop above] node {0-9} ()
               edge[bend left=20] node {\kaUS} (UInt)
               edge node {.}           (NumDot)
               edge[bend left=8] node[swap, pos=0.72] {e,E} (Exp)
      (UInt)   edge[loop right] node {\kaUS} ()
               edge[bend left=20] node {0-9} (Int)
      % desetinná část
      (NumDot) edge node {0-9}         (Double)
      (Period) edge node[swap] {0-9}   (Double)
      (Double) edge[out=-60, in=-20, loop] node[below right] {0-9} ()
               edge[bend left=15] node {\kaUS} (UDouble)
               edge node[swap] {e,E}   (Exp)
               edge node {d,D}         (DFinal)
      (UDouble) edge[loop right] node {\kaUS} ()
               edge[bend left=15] node {0-9} (Double)
      % exponent
      (Exp)    edge node {+,-}         (SExp)
               edge node {0-9}         (EDouble)
      (SExp)   edge node {0-9}         (EDouble)
      (EDouble) edge[loop right] node {0-9} ()
               edge[bend left=15] node {\kaUS} (EUDouble)
               edge node {d,D}         (DFinal)
      (EUDouble) edge[loop above] node {\kaUS} ()
               edge[bend left=15] node {0-9} (EDouble);
  \end{tikzpicture}
  \caption{KA -- část 1: číselné literály (\texttt{Int}, \texttt{Double}) a tečka.}
  \label{fig:ka}
\end{figure}

\begin{figure}[H]
  \centering
  \begin{tikzpicture}[ka]
    \node[state, initial]   (S)       at (0,0)      {S};
    \node[state]            (Q1)      at (2.2,0)    {Q1};
    \node[state]            (In)      at (5,1.8)    {In};
    \node[state]            (Esc)     at (5,4)      {Esc};
    \node[state, accepting] (StrN)    at (8.2,1.8)  {StrNormal};
    \node[state, accepting] (StrE)    at (4.4,-2)   {StrEmpty};
    \node[state]            (InM)     at (7.4,-2)   {InMultiline};
    \node[state]            (Q1M)     at (10.4,-2)  {Q1Multi};
    \node[state]            (Q2M)     at (13.2,-2)  {Q2Multi};
    \node[state, accepting] (StrM)    at (13.2,0.4) {StrMultiline};

    \path
      (S)    edge node {\kaQ}             (Q1)
      (Q1)   edge node {C}                (In)
             edge node {\kaBS}            (Esc)
             edge node[swap] {\kaQ}       (StrE)
      (In)   edge[loop below] node {C}    ()
             edge[bend left=20] node {\kaBS} (Esc)
             edge node {\kaQ}             (StrN)
      (Esc)  edge[bend left=20] node {\kaBS, \kaQ, n, r, t} (In)
      (StrE) edge node {\kaQ}             (InM)
      (InM)  edge[loop above] node {M}    ()
             edge node {\kaQ}             (Q1M)
      (Q1M)  edge node {\kaQ}             (Q2M)
             edge[bend left=30] node {M}  (InM)
      (Q2M)  edge node[swap] {\kaQ}       (StrM)
             edge[bend left=40] node {M}  (InM)
      (StrM) edge[loop above] node {\kaQ} ();
  \end{tikzpicture}
  \caption{KA -- část 2: řetězcové literály (jednořádkové a víceřádkové).}
  \label{fig:ka-str}
\end{figure}

\begin{figure}[H]
  \centering
  \begin{tikzpicture}[ka]
    \node[state, initial]   (S)      at (0,0)      {S};
    \node[state, accepting] (Und)    at (2.8,1.8)  {Underscore};
    \node[state, accepting] (I)      at (6,0)      {I};
    \node[state]            (At)     at (2.4,-2)   {Q@};
    \node[state]            (Atm)    at (4.6,-2)   {Q@m};
    \node[state]            (Atma)   at (6.8,-2)   {Q@ma};
    \node[state]            (Atmai)  at (9,-2)     {Q@mai};
    \node[state, accepting] (Atmain) at (11.4,-2)  {Q@main};

    \path
      (S)     edge node {\kaUS}             (Und)
              edge node {A-Z, a-z}       (I)
              edge node[swap] {@}        (At)
      (Und)   edge node {A-Z, a-z, 0-9, \kaUS} (I)
      (I)     edge[loop right] node {A-Z, a-z, 0-9, \kaUS} ()
      (At)    edge node {m}              (Atm)
      (Atm)   edge node {a}              (Atma)
      (Atma)  edge node {i}              (Atmai)
      (Atmai) edge node {n}              (Atmain);
  \end{tikzpicture}
  \caption{KA -- část 3: identifikátory, zástupný symbol \texttt{\_} a anotace \texttt{@main}.}
  \label{fig:ka-id}
\end{figure}

\begin{figure}[H]
  \centering
  \begin{tikzpicture}[ka]
    \node[state, initial]   (S)    at (0,0)     {S};
    \node[state, accepting] (Eq)   at (3,3)     {Q=};
    \node[state, accepting] (EqEq) at (6,3.6)   {Q==};
    \node[state, accepting] (Arr)  at (6,2.4)   {Q=>};
    \node[state, accepting] (Not)  at (3,1.4)   {Q!};
    \node[state, accepting] (Neq)  at (6,1.4)   {Q!=};
    \node[state, accepting] (Gt)   at (3,0.2)   {Q>};
    \node[state, accepting] (Ge)   at (6,0.2)   {Q>=};
    \node[state, accepting] (Lt)   at (3,-1.4)  {Q<};
    \node[state, accepting] (Le)   at (6,-0.8)  {Q<=};
    \node[state, accepting] (Gen)  at (6,-2)    {Q<-};
    \node[state]            (Amp)  at (3,-3.2)  {Q\&};
    \node[state, accepting] (And)  at (6,-3.2)  {Q\&\&};
    \node[state, accepting] (Bar)  at (3,-4.4)  {Q|};
    \node[state, accepting] (Or)   at (6,-4.4)  {Q||};

    \path
      (S)   edge node {=}        (Eq)
            edge node[pos=0.6] {!}  (Not)
            edge node[pos=0.6] {>}  (Gt)
            edge node[pos=0.6] {<}  (Lt)
            edge node[pos=0.6] {\&} (Amp)
            edge node[swap] {|}     (Bar)
      (Eq)  edge node {=}        (EqEq)
            edge node[swap] {>}  (Arr)
      (Not) edge node {=}        (Neq)
      (Gt)  edge node {=}        (Ge)
      (Lt)  edge node {=}        (Le)
            edge node[swap] {-}  (Gen)
      (Amp) edge node {\&}       (And)
      (Bar) edge node {|}        (Or);
  \end{tikzpicture}
  \caption{KA -- část 4: relační, logické a další víceznakové operátory.}
  \label{fig:ka-op}
\end{figure}

\begin{figure}[H]
  \centering
  \begin{tikzpicture}[ka]
    \node[state, initial above] (S) at (0,0) {S};
    \foreach \name/\lab/\ang in {
        Plus/Q+/0, Minus/Q-/30, Mul/Q*/60, Div/{Q/}/120, LPar/Q(/150,
        RPar/Q)/180, LBr/Q[/210, RBr/Q]/240, Colon/Q:/270, Semi/Q;/300,
        Comma/Comma/330} {
      \node[state, accepting] (\name) at (\ang:3) {\lab};
    }
    \node[state, accepting] (Concat) at (0:5.6) {Q++};

    \path
      (S)    edge node[swap, pos=0.6] {+} (Plus)
             edge node[swap, pos=0.6] {-} (Minus)
             edge node[swap, pos=0.6] {*} (Mul)
             edge node[swap, pos=0.6] {/} (Div)
             edge node[swap, pos=0.6] {(} (LPar)
             edge node[swap, pos=0.6] {)} (RPar)
             edge node[swap, pos=0.6] {[} (LBr)
             edge node[swap, pos=0.6] {]} (RBr)
             edge node[swap, pos=0.6] {:} (Colon)
             edge node[swap, pos=0.6] {;} (Semi)
             edge node[swap, pos=0.6] {,} (Comma)
      (Plus) edge node {+} (Concat);
  \end{tikzpicture}
  \caption{KA -- část 5: aritmetické operátory, závorky a oddělovače.}
  \label{fig:ka-delim}
\end{figure}

\begin{figure}[H]
  \centering
  \begin{tikzpicture}[ka]
    \node[state, initial]   (S)   at (0,0)    {S};
    \node[state, accepting] (NL)  at (3,0)    {NewLine};
    \node[state, accepting] (NLS) at (7.5,1.6)  {NewLineSpaces};
    \node[state]            (NLT) at (7.5,-1.6) {NewLineTabs};

    \path
      (S)   edge node {\kaLF}                   (NL)
      (NL)  edge[loop above] node {\kaLF}       ()
            edge[bend left=15] node {\kaSP}     (NLS)
            edge[bend right=15] node[swap] {\kaTAB} (NLT)
      (NLS) edge[loop right] node {\kaSP}       ()
            edge[bend left=15] node[pos=0.25] {\kaLF} (NL)
            edge node {\kaTAB}                  (NLT)
      (NLT) edge[loop right] node {\kaSP, \kaTAB} ()
            edge[bend right=15] node[swap, pos=0.25] {\kaLF} (NL);
  \end{tikzpicture}
  \caption{KA -- část 6: konec řádku a odsazení na začátku řádku.}
  \label{fig:ka-nl}
\end{figure}


%------------------------------------------------------------
\section{Syntaktická analýza}

Syntaktickou analýzu jsme implementovali metodou rekurzivního sestupu. Pro každý neterminál gramatiky existuje jedna
funkce (\texttt{p\_program}, \texttt{p\_function}, \texttt{p\_code\_block},
\texttt{p\_statement}, \texttt{p\_if}, \texttt{p\_while}, \texttt{p\_for},
\ldots) a o~dalším postupu rozhoduje vždy jeden token dopředu, který drží
lexer. Bloky kódu jsou dané odsazením, které lexer převádí na tokeny
\texttt{INDENT} a \texttt{DEDENT}, takže je parser zpracovává jako běžné
závorky. Během analýzy parser rovnou staví abstraktní syntaktický strom
a plní tabulky symbolů. Funkci lze volat i před její definicí: volání si
v~tabulce funkcí založí záznam, který definice později doplní.

Výrazy zpracovává funkce \texttt{p\_expression} pomocí precedenční
syntaktické analýzy zdola nahoru. Na zásobníku terminálů je na dně
zarážka \texttt{\$}. O~tom, zda se operátor ze vstupu přesune na zásobník
(\emph{take}), nebo se má redukovat (\emph{reduce}), rozhoduje precedenční
tabulka $13\times13$ pro operátory \texttt{*}, \texttt{/}, \texttt{+},
\texttt{-}, \texttt{++}, relační operátory, \texttt{==}/\texttt{!=},
\texttt{\&\&}, \texttt{||}, \texttt{to}, \texttt{until}, \texttt{by}
a \texttt{\$}. Relační operátory a operátory rovnosti jsou neasociativní,
takže například \texttt{a < b < c} skončí syntaktickou chybou. Každá
redukce vytvoří uzel binární operace v~AST\@.

Operandy v~tabulce nejsou. Analyzátor střídá dva stavy, „čekám operand“
a „čekám operátor“, a ve stavu operandu volá rekurzivní sestup
(\texttt{p\_operand}). Ten zpracuje literály, proměnné, volání funkcí,
indexování, výrazy v~závorkách i unární operátory. Díky tomu se unární
a binární minus rozliší podle stavu, ne podle tokenu. Všechny tokeny, které do výrazu nepatří (nový řádek, \texttt{then}, \texttt{do}, čárka, pravá závorka
apod.) se chovájí jako konec vstupu (\texttt{\$}), zbytek zásobníku se
zredukuje a řízení se vrátí rekurzivnímu sestupu.

\subsection{LL gramatika}
LL gramatiku jsme nepsali dopředu jako samostatný dokument, ale rovnou ji
převáděli do kódu rekurzivního sestupu.

\subsection{LL tabulka}
LL tabulku jsme nepoužili, protože jsme parser implementovali s metodou rekurzivního sestupu.

\subsection{Precedenční tabulka}
\label{sec:prec}
\begin{table}[H]
  \centering
  \setlength{\tabcolsep}{4pt}
  \begin{tabular}{l|*{13}{c}}
    \toprule
    Zásobník \textbackslash Vstup & \texttt{*} & \texttt{/} & \texttt{+} & \texttt{-} & \texttt{++} & \texttt{<=} & \texttt{==}
    & \texttt{\&\&} & \texttt{||} & \texttt{to} & \texttt{until} & \texttt{by} & \$ \\
    \midrule
    \texttt{*}     & > & > & > & > & > & > & > & > & > & > & > & > & > \\
    \texttt{/}     & > & > & > & > & > & > & > & > & > & > & > & > & > \\
    \texttt{+}     & < & < & > & > & > & > & > & > & > & > & > & > & > \\
    \texttt{-}     & < & < & > & > & > & > & > & > & > & > & > & > & > \\
    \texttt{++}    & < & < & > & > & > & > & > & > & > & > & > & > & > \\
    \texttt{<=}    & < & < & < & < & < & -- & > & > & > & > & > & > & > \\
    \texttt{==}    & < & < & < & < & < & < & -- & > & > & > & > & > & > \\
    \texttt{\&\&}  & < & < & < & < & < & < & < & > & > & > & > & > & > \\
    \texttt{||}    & < & < & < & < & < & < & < & < & > & > & > & > & > \\
    \texttt{to}    & < & < & < & < & < & < & < & < & < & > & > & > & > \\
    \texttt{until} & < & < & < & < & < & < & < & < & < & < & > & > & > \\
    \texttt{by}    & < & < & < & < & < & < & < & < & < & < & < & > & > \\
    \$             & < & < & < & < & < & < & < & < & < & < & < & < & -- \\
    \bottomrule
  \end{tabular}
  \caption{Precedenční tabulka.
           < = načtení, > = redukce,
           -- = chyba. V záhlavích řádků i sloupců \texttt{<=} zastupuje operátory \texttt{<}, \texttt{<=}, \texttt{>} a \texttt{>=}. \texttt{==} zastupuje i operátor \texttt{!=}.}
  \label{tab:prec}
\end{table}

\section{Sémantická analýza}
\label{sec:semantika}
Sémantickou analýzu provádíme ve dvou fázích. Kontroly, které lze udělat
hned při čtení programu, dělá přímo parser (redefinice a použití před definicí). Zbytek běží jako samostatný
průchod nad hotovým AST (\texttt{semantic.c}), protože k~němu je potřeba
znát signatury všech funkcí. Funkce totiž lze volat ještě před jejich
definicí.

\subsection{Samostatný sémantický průchod}
Po úspěšném parsování projde funkce \texttt{semantic\_check} program
v~několika krocích:

\begin{enumerate}
  \item Doplní vestavěné funkce do tabulky funkcí.
  \item Ověří, že každá volaná funkce je definovaná a že existuje
        \texttt{@main}.
  \item Z~AST sestaví graf volání (funkce $\rightarrow$ funkce, které volá)
        a poznamená si, které funkce obsahují \texttt{return}.
  \item Funkce bez explicitního návratového typu nesmí obsahovat
        \texttt{return} ani ležet na cyklu grafu volání, tedy nesmí být
        přímo ani nepřímo rekurzivní. Cyklus hledáme obyčejným
        průchodem do hloubky z~dané funkce.
  \item Rekurzivně typově zkontroluje těla všech funkcí a odvodí chybějící návratové typy. Máme zkontrolováno, že v závislosti typů nejsou smyčky, takže nedojde k zacyklení.
\end{enumerate}

\subsection{Typová kontrola}
Typy výrazů počítáme zdola nahoru a ukládáme je přímo do uzlů AST, odkud je
pak přebírá generátor kódu. Proměnná definovaná bez anotace
(\texttt{val x = e}) převezme typ inicializačního výrazu. U implicitních typových konverzí doplníme do AST příslušný \texttt{ASTNode}.

Do funkcí bez explicitního returnu zkontrolujeme pravidla a doplníme implicitní \texttt{return}, popřípadě \texttt{return Unit}.

\section{Generátor cílového kódu}
Cílový kód v jazyce IFJcode26 je generován přímo ze struktury AST a tabulky symbolů bez použití vnitřního mezikódu. Není prováděna optimalizace pro velikost či ryclost interpretace.

Generace je realizována rekurzivním průchodem AST. Pro vyhodnocování výrazů, předávaní parametů a návratových hodnot funkcí je využit datový zásobník. Jelikož je cílový kód generován přímo do výsledného souboru, bylo by vytváření dočasných proměnných složité a nepraktické. Proto jsou pro veškeré operace bez podpory datového zásobníku využity globální proměnné interpretu, sloužící jako dočasné proměnné. Jsou využity především pro práci s listy. 

%------------------------------------------------------------
\section{Popis modulů}
\label{sec:moduly}
% Informačně obohacující věta na dvojici .c/.h.
\begin{tabularx}{\textwidth}{@{}l X@{}}
    \toprule
    \textbf{Modul} & \textbf{Popis}\\
    \midrule
    \texttt{main.c} & Vstupní bod programu.\\
    \texttt{ast.c/.h} & Definice a operace nad abstraktním syntaktickým stromem.\\
    \texttt{bst.c/.h} & Implementace výškově vyváženého binárního vyhledávacího stromu.\\
    \texttt{impl\_bst.h} & Interní definice uzlů a struktury binárního vyhledávacího stromu.\\
    \texttt{builtins.c/.h} & Signatury vestavěných funkcí a jejich cílový kód.\\
    \texttt{codegen.c/.h} & Generování cílového kódu z abstraktního syntaktického stromu.\\
    \texttt{codegen\_globals.c/.h} & Globální proměnné cílového kódu používané jako dočasné proměnné.\\
    \texttt{dynamic\_array.c/.h} & Správa kapacity generického dynamického pole.\\
    \texttt{error\_trace.c/.h} & Výpis původu chyby a řetězce jejího šíření.\\
    \texttt{function\_table.c/.h} & Operace nad tabulkou funkcí programu.\\
    \texttt{instruction.c/instructions.h} & Reprezentace operandů a výpis instrukcí cílového kódu.\\
    \texttt{keyword.c/.h} & Tabulka klíčových slov a jejich rozpoznávání lexikálním analyzátorem.\\
    \texttt{lexer.c/.h} & Lexikální analýza zdrojového kódu IFJ26.\\
    \texttt{lexer\_dfa.c/.h} & Tabulky stavů a přechodů konečného automatu lexikálního analyzátoru.\\
    \texttt{parser.c/.h} & Syntaktická analýza zdrojového kódu IFJ26.\\
    \texttt{psa.c/.h} & Zásobník precedenční analýzy a redukce výrazů.\\
    \texttt{semantic.c/.h} & Sémantická analýza, kontrola typů a pravidel volání funkcí.\\
    \texttt{source\_preprocess.c/.h} & Předzpracování zdrojového kódu, kontrola importů a odstranění komentářů.\\
    \texttt{symtable.c/.h} & Operace nad tabulkou symbolů.\\
    \texttt{type.c/.h} & Definice sémantických typů a řešení jejich závislostí.\\
    \bottomrule
\end{tabularx}

%------------------------------------------------------------
\section{Využití prostředků AI}
\label{sec:ai}
AI jsme používali hlavně jako agentního pomocníka při implementaci a při
konzultacích nad architekturou systému. Dále jsme ji používali pro code-reviews a kontroly návrhů architektury. Hodně nám pomohla i při ujasňování
nejasných míst ve specifikaci, kde jsme použili nástroj NotebookLM.

Největší přínos měla AI při psaní testů. Repozitář obsahuje 467 jednotkových
testů v~GoogleTestu a 214 end-to-end případů (programů v~IFJ26 s~očekávaným
výstupem nebo návratovým kódem), které pokrývají lexikální, syntaktickou
i sémantickou analýzu, vestavěné funkce a celé programy. Většinu z~nich
vygenerovala AI podle našich zadání. Zkušenost s~AI hodnotíme celkově
pozitivně, výrazně nám zrychlila hlavně implementaci a testování.

\subsection{Použité nástroje a modely}
\begin{itemize}
  \item Claude Sonnet 5.5 ve webovém rozhraní (konzultace návrhu),
  \item Google Gemini notebook (dříve NotebookLM), model neuveden
        (dotazy nad zadáním projektu),
  \item Google Antigravity, model Gemini 3.8 Flash,
  \item Opencode, model Kimi-k3.1-pro,
  \item Claude Code, model Claude Opus 5.5 (implementace a naprostá většina
        testů),
  \item OpenAI Codex přes Hermes Agent, model gpt-5.6-terra (část testů
        lexeru),
  \item Claude Fable 5.2.
\end{itemize}

\subsection{Typický pracovní postup}
Práce s~AI u~nás probíhala v~několika ustálených krocích:

\begin{enumerate}
  \item \textbf{Ujasnění zadání.} Když nám nebylo jasné, co specifikace
        přesně vyžaduje, zeptali jsme se nástroje NotebookLM, do kterého jsme
        nahráli zadání projektu. Jeho odpovědi nám sloužily jen jako vodítko
        a vždy jsme je ověřili přímo v~textu specifikace.
  \item \textbf{Návrh v~týmu.} Architekturu a rozhraní modulů jsme navrhovali
        sami na týmových poradách. AI nám přitom sloužila spíš jako partner do
        diskuze, se kterým jsme procházeli různé varianty řešení. Například
        konečný automat lexeru jsme nejdřív formálně popsali v~samostatném
        dokumentu (\texttt{KA.md}) a lexer pak podle tohoto popisu
        implementovali jako tabulkou řízený automat.
  \item \textbf{Malé, jasně vymezené úkoly.} Agentovi (Claude Code, Opencode,
        Antigravity) jsme nikdy nezadávali „napiš parser“. Dostal vždy menší
        úkol a k~němu potřebný kontext: příslušnou část specifikace, hlavičkové
        soubory s~rozhraním a ukázku stávajícího kódu, aby výsledek odpovídal
        stylu projektu.
  \item \textbf{Ověření.} Vygenerovaný kód musel odpovídat očekávání zadavatele.          Dále se musel dát přeložit a projít testy.
        Formátování (\texttt{clang-format}) a statickou analýzu
        (\texttt{cppcheck}) kontrolovaly hooky \texttt{pre-commit}, které
        se spouštěly lokálně i v~GitHub Actions při každém pushi.
  \item \textbf{Code review.} Do hlavní větve se kód od AI dostával stejně
        jako každý jiný, tedy přes pull request a kontrolu dalším členem týmu.
  \item \textbf{Upřesnění místo záplatování.} Když výsledek neodpovídal naší
        představě, většinou jsme ho neopravovali po kouscích. Zpřesnili jsme
        zadání a nechali kód vygenerovat znovu.
\end{enumerate}

U~testů jsme postupovali podobně, jen jsme očekávané výsledky vždy odvozovali
ze specifikace, nikdy z~toho, jak se překladač zrovna choval. Díky tomu testy
odhalily několik míst, kde se implementace od specifikace odchylovala,
například nerozpoznané komentáře, chybné zpracování čísel s~úvodní nulou nebo
přijetí literálu typu \texttt{Double} mimo povolený rozsah. Takové testy jsme
ponechali s~původním očekáváním a do opravy je dočasně vypnuli (prefix
\texttt{DISABLED\_}). Chybu pak opravil člen týmu, který měl daný modul
na starosti. Testy jsme přidávali po menších tematických dávkách, abychom
je dokázali zkontrolovat. Každý test vygenerovaný AI má v~hlavičce uvedeno,
jaký model ho vytvořil a s~jakým zadáním.

Obecně platilo, že AI nám výrazně zrychlila psaní kódu, ale rozhodnutí
o~návrhu a odpovědnost za výsledek zůstaly na nás.

\subsection{Problémy při využití AI}
Rozsáhlejší úsek kódu vygenerovaný AI je potřeba následně zkontrolovat, což
bývá náročné a demotivující. Vygenerovaný kód je navíc podle naší zkušenosti
skoro vždy na první pohled správný, a tak má recenzent sklon kontrolu
odbývat a přehlédnout chybu. To platí hlavně ve složitém systému s~mnoha
mezními případy, jakým překladač je.

Podobné riziko jsme viděli u~testů. Pokud by model dostal k~dispozici jen
zdrojový kód, snadno by očekávané výsledky opsal ze stávajícího chování,
a test by pak chybu nenašel, ale „zakonzervoval“. Proto jsme modelu vždy
dávali příslušnou část specifikace a výsledky odvozovali z~ní. Obecně se
ukázalo, že klíčové je mít specifikaci jasně sepsanou a poctivě ji ověřovat
testy.

Odpovědi nad zadáním (NotebookLM) byly užitečné, ale ne vždy přesné, a proto
jsme je nikdy nebrali jako jediný zdroj pravdy.

\subsection{Návrhy AI nevyužité nebo výrazně upravené}
Původní implementace lexeru vygenerovaná AI byla velmi nepřehledná
a neodpovídala naší představě, takže jsme ji nepoužili. Místo toho jsme
nejdřív formálně popsali konečný automat (stavy, abecedu a přechody
v~dokumentu \texttt{KA.md}) a rozhodli jsme se, že bude v~kódu
reprezentován tabulkou přechodů. Teprve podle tohoto přesnějšího zadání
vznikla implementace, kterou projekt používá.

U~testů, které neprocházely, AI často navrhovala „upravit očekávaný
výsledek podle překladače“. Takové řešení jsme vědomě nepřijímali:
pokud test odpovídal specifikaci, opravovali jsme implementaci, ne test.

Dokumenty \texttt{specifikace.txt} a \texttt{psa\_spec.md} jsme
sepisovali s~pomocí AI agentů. Sloužily nám jako podrobné návrhy
implementace datových struktur a jednotlivých částí překladače.
Při jejich vypracování jsme spolu s~agenty často narazili na logické
mezery, takže oba dokumenty před samotnou implementací prošly mnoha
úpravami.

%------------------------------------------------------------
\end{document}